% Independent preview; original geometry and numeric relationships retained.
% 原创计数示例：二维单元含3点，其x1投影含5点，阈值tau=3。
\begin{tikzpicture}[x=12mm,y=10mm,font=\small,text=FigureInk,
  line width=0.7pt]
 \fill[FigureYellow!12] (1,1) rectangle (2,2);
 \draw[step=1,FigureGrid] (0,0) grid (4,3);
 \draw[LancetGreen,thick] (1,1) rectangle (2,2);
 \draw[->,FigureMuted] (0,0) -- (4.5,0) node[right] {$x_1$};
 \draw[->,FigureMuted] (0,0) -- (0,3.4) node[above] {$x_2$};
 \foreach \x/\y in {1.25/1.25,1.7/1.4,1.4/1.8}
  \fill[LancetGreen] (\x,\y) circle[radius=2pt];
 \foreach \x/\y in {1.2/0.45,1.8/2.55,0.4/1.3,3.1/2.3}
  \draw[FigureBlue,fill=white] (\x,\y) circle[radius=2pt];
 \draw[FigureMuted,dashed] (1,0) -- (1,-1.2) (2,0) -- (2,-1.2);
 \draw[FigureMuted] (0,-1.2) -- (4,-1.2);
 \draw[LancetGreen,line width=2pt] (1,-1.2) -- (2,-1.2);
 \foreach \x in {1.2,1.25,1.4,1.7,1.8}
  \fill[FigureInk] (\x,-1.2) circle[radius=1.6pt];
 \node[right,align=left] at (4.6,1.65) {二维单元\\计数 $3$};
 \node[right,align=left] at (4.6,-1.2) {一维投影\\计数 $5$};
 \draw[book arrow] (3,-0.25) -- (3,-0.95)
   node[right,midway] {去掉 $x_2$ 限制};
\end{tikzpicture}
